\section{Related Work}

\subsection{Zero-Shot Learning in Fault Diagnosis}

Zero-Shot Learning (ZSL) has emerged as a promising solution to address the paucity of annotated fault data in industrial environments, where acquiring samples of rare or catastrophic faults is often costly or prohibited. The core objective of ZSL in condition monitoring is to establish a transferable knowledge bridge between seen categories and unseen categories. This transfer is facilitated through an intermediate semantic space \cite{mikolov2013efficient,oord2018representation}, typically comprising human-defined attributes or linguistic descriptors that encapsulate the physical characteristics of the faults \cite{zhang2022intelligent}.

Early explorations into zero-shot fault diagnosis relied on direct nearest-neighbor searches within a shared embedding space, associating query signals with the closest unseen class prototypes via cosine similarity or Euclidean distance \cite{tzeng2017adversarial}. However, these projection-based techniques frequently suffer from the "hubness" problem \cite{xian2019zero} and severe domain shift \cite{tzeng2017adversarial} when dealing with the non-stationary, non-linear signals \cite{li2024clip} inherent to rotating machinery.

To mitigate these issues, generative ZSL models \cite{xian2018feature} have gained traction. By synthesizing pseudo-samples for unseen classes conditioned on semantic attributes, these models convert the ZSL problem into a conventional supervised task \cite{liao2024cyclegan}. Variational Autoencoders (VAEs) \cite{lv2020hybrid} and Generative Adversarial Networks (GANs) are central to this paradigm. For instance, attribute-based VAEs learn latent distributions to generate unseen fault features \cite{wang2020deep}, while frameworks like CycleGAN-SD \cite{liao2024cyclegan,zhu2017unpaired} enforce cycle-consistency between signal and semantic \cite{radford2021learning} spaces to preserve the structural integrity of time-series data.

Despite their success, GAN and VAE-based ZSL methodologies face two fundamental bottlenecks. First, reliance on manually engineered discrete attribute matrices restricts the expressiveness of the semantic space and introduces subjective bias. Second, GANs are prone to mode collapse and optimization instability \cite{arjovsky2017wasserstein}, which are exacerbated when synthesizing high-frequency industrial signals. The adversarial min-max game fails to guarantee sufficient diversity, leading to generated distributions that may lack physical plausibility \cite{goodfellow2014generative}. Consequently, a more stable generative paradigm coupled with a continuous, expressive semantic space is required.

\subsection{Diffusion Models for Condition Monitoring}

Denoising Diffusion Probabilistic Models (DDPMs) have revolutionized deep generative modeling, surpassing GANs in synthesis quality, distributional coverage, and training stability \cite{ho2020denoising}. Diffusion models learn to progressively inject Gaussian noise into data via a forward Markov chain and subsequently learn a parameterized reverse denoising process to reconstruct the original distribution from pure noise. This transition from adversarial training to likelihood-based diffusion processes represents a significant leap forward for time-series analysis and condition monitoring.

In the forward process, a real fault signal is gradually transformed into Gaussian noise over a sequence of steps according to a variance schedule. The generative reverse process then employs a neural network (typically a U-Net adapted for 1D sequences) to predict and subtract the noise at each step, effectively reconstructing the signal from an isotropic Gaussian state \cite{he2016deep,ho2020denoising}. This training objective provides a more stable optimization landscape than adversarial equivalents and ensures better coverage of the data distribution.

Recent literature has begun exploring DDPMs for fault data augmentation, demonstrating that they can generate high-fidelity, physically consistent vibration signals that respect underlying mechanical kinematics better than GANs \cite{fan2025lightweight,yang2024novel}. However, the application of DDPMs to zero-shot learning remains largely unaddressed. Unconditional diffusion cannot synthesize specific unseen categories on demand, and existing conditional models primarily utilize discrete class labels. Bridging the gap between continuous linguistic descriptions and the reverse denoising path is non-trivial. Moreover, current time-series diffusion architectures often focus on the time domain, neglecting the discriminative frequency signatures crucial for diagnosing rotating machinery faults. This underscores the need for a conditional diffusion framework tailored for cross-modal zero-shot synthesis that captures dual-path dynamics.

\subsection{Vision-Language Models in Time-Series Analysis}

Large-scale Vision-Language Models (VLMs), such as CLIP, have demonstrated the potential of aligning distinct modalities within a unified semantic embedding space \cite{mikolov2013efficient}. By leveraging massive corpora of paired data, these architectures learn continuous, expressive representations that facilitate zero-shot classification without task-specific fine-tuning \cite{radford2021learning}. This paradigm shift has catalyzed advancements in industrial signal processing, recognizing that Large Language Models (LLMs) \cite{brown2020language,devlin2019bert} possess innate knowledge of physical systems and fault mechanisms acquired during technical pre-training \cite{liu2021generalized,liu2023dsd}.

By querying an LLM with structured prompts detailing fault characteristics and root causes \cite{shi2024unseen,luo2024self,qin2024feature}, one can extract dense, continuous textual embeddings that are more informative than traditional discrete binary matrices \cite{brown2020language}. In multi-modal diagnostics, the objective is to align these textual embeddings with abstract signal representations. This transition shifts ZSL from rigid attribute regression toward flexible cross-modal contrastive learning \cite{chen2020simple}. Typically, signal-to-text and text-to-signal alignments are optimized via symmetric contrastive loss functions, such as InfoNCE, which maximize the similarity of paired embeddings while penalizing hard negatives in the latent space \cite{oord2018representation,chen2020simple}.

While contrastive techniques are effective for discriminative tasks like anomaly detection, their integration into generative ZSL for signal synthesis is largely unexplored. Injecting continuous, contrastively-aligned semantic embeddings as conditioning vectors into the reverse path of a diffusion model presents a novel mechanism to guide the synthesis of unseen fault signals. This approach bypasses the limitations of discrete attributes and forces the generator to respect high-level textual semantics parsed by the LLM. Furthermore, employing a contrastive consistency constraint on generated signals ensures they preserve their semantic anchors, resolving domain shift issues prevalent in conventional ZSL. The proposed CDDM framework weaves these concepts together, establishing an innovative text-to-signal generative pathway for complex industrial monitoring scenarios.